\documentclass[DIV=calc]{scrartcl}
\usepackage[utf8]{inputenc}
\usepackage[T1]{fontenc}
\usepackage[ngerman]{babel}
\usepackage{graphicx}
\usepackage[draft, markup=underlined]{changes}
\usepackage{csquotes}
\usepackage{eurosym}

\usepackage{ulem}
%\usepackage[dvipsnames]{xcolor}
\usepackage{paralist}
%\usepackage{fixltx2e}
%\usepackage{ellipsis}
\usepackage[tracking=true]{microtype}

\usepackage{lmodern}              % Ersatz fuer Computer Modern-Schriften
%\usepackage{hfoldsty}

%\usepackage{fourier}             % Schriftart
\usepackage[scaled=0.81]{helvet}     % Schriftart

\usepackage{url}
%\usepackage{tocloft}             % Paket für Table of Contents
\def\UrlBreaks{\do\a\do\b\do\c\do\d\do\e\do\f\do\g\do\h\do\i\do\j\do\k\do\l%
\do\m\do\n\do\o\do\p\do\q\do\r\do\s\do\t\do\u\do\v\do\w\do\x\do\y\do\z\do\0%
\do\1\do\2\do\3\do\4\do\5\do\6\do\7\do\8\do\9\do\-}%

\usepackage{xcolor}
\definecolor{urlred}{HTML}{660000}

\usepackage{hyperref}
\hypersetup{colorlinks=false}

%\usepackage{mdwlist}     % Änderung der Zeilenabstände bei itemize und enumerate
% \usepackage{draftwatermark} % Wasserzeichen ``Entwurf''
% \SetWatermarkText{Antrag}

\parindent 0pt                 % Absatzeinrücken verhindern
\parskip 12pt                 % Absätze durch Lücke trennen

\setlength{\textheight}{23cm}
\usepackage{fancyhdr}
\pagestyle{fancy}
\fancyhead{} % clear all header fields
\cfoot{}
\lfoot{Zusammenkunft aller Physik-Fachschaften}
\rfoot{www.zapfev.de\\stapf@zapf.in}
\renewcommand{\headrulewidth}{0pt}
\renewcommand{\footrulewidth}{0.1pt}
\newcommand{\gen}{*innen}
\addto{\captionsngerman}{\renewcommand{\refname}{Quellen}}

%%%% Mit-TeXen Kommandoset
\usepackage[normalem]{ulem}
\usepackage{xcolor}
\usepackage{xspace} 

\newcommand{\replace}[2]{
    \sout{\textcolor{blue}{#1}}~\textcolor{blue}{#2}}
\newcommand{\delete}[1]{
    \sout{\textcolor{red}{#1}}}
\newcommand{\add}[1]{
    \textcolor{red}{#1}}

\newif\ifcomments
\commentsfalse
%\commentstrue

\newcommand{\red}[1]{{\ifcomments\color{red} {#1}\else{#1}\fi}\xspace}
\newcommand{\blue}[1]{{\ifcomments\color{blue} {#1}\else{#1}\fi}\xspace}
\newcommand{\green}[1]{{\ifcomments\color{green} {#1}\else{#1}\fi}\xspace}

\newcommand{\repl}[2]{{\ifcomments{\color{red} \sout{#1}}{\color{blue} {\xspace #2}}\else{#2}\fi}}
%\newcommand{\repl}[2]{{\color{red} \sout{#1}\xspace{\color{blue} {#2}}\else{#2}\fi}\xspace}

\newcommand{\initcomment}[2]{%
	\expandafter\newcommand\csname#1\endcsname{%
		\def\thiscommentname{#1}%
		\definecolor{col}{rgb}{#2}%
		\def\thiscommentcolor{col}%
}}

% initcomment Name RGB-color
\initcomment{Philipp}{0, 0.5, 0}

%\renewcommand{\comment}[1]{{\ifcomments{\color{red} {#1}}{}\fi}\xspace}

\renewcommand{\comment}[2][\nobody]{
	\ifdefined#1
	{\ifcomments{#1 \expandafter\color{\thiscommentcolor}{\thiscommentname: #2}}{}\fi}\xspace
	\else
	{\ifcomments{\color{red} {#2}}{}\fi}\xspace
	\fi
}

\newcommand{\zapf}{ZaPF\xspace}

\let\oldgrqq=\grqq
\def\grqq{\oldgrqq\xspace}

\setlength{\parskip}{.6em}
\setlength{\parindent}{0mm}

%\usepackage{geometry}
%\geometry{left=2.5cm, right=2.5cm, top=2.5cm, bottom=3.5cm}

% \renewcommand{\familydefault}{\sfdefault}




\begin{document}

\hspace{0.87\textwidth}
\begin{minipage}{120pt}
	\vspace{-1.8cm}
	\includegraphics[width=80pt]{logo.pdf}
	\centering
	\small Zusammenkunft aller Physik-Fachschaften
\end{minipage}

\begin{center}
  \huge{Studierende in der Krise nicht vergessen!}\vspace{.25\baselineskip}\\
  \normalsize
\end{center}
\vspace{1cm}

%%%% Metadaten %%%%

% \paragraph{Adressierte:} Ampelkoalition, deren Bildungspoltische Sprecher*innen, BMBF, fzs, MeTaFa


% \paragraph{Antragstellende:} Christian (UMR), Vale (Köln), Kai (JLU Gießen), Lennart (Heidelberg), Noemi (Heidelberg)

% %%%% Text des Antrages zur veröffentlichung %%%%

% \section*{Antragstext}

Die ZaPF fordert, dass die Bedürfnisse, Sorgen und finanziellen Nöte der Studierenden gehört werden und daraufhin adäquat gehandelt wird. Nach über drei Jahren Coronapandemie mit vielen Einschränkungen fühlen wir uns nicht wahrgenommen\footnote{\url{https://katapult-magazin.de/de/artikel/was-die-medien-verschwiegen-haben}}.
%{https://katapult-magazin.de/de/artikel/was-die-medien-verschwiegen-haben}}. 
Vielmehr hat die damalige Bundesregierung die Auszahlung von Coronahilfen für Studierende unnötig hinausgezögert\footnote{\url{https://www.zeit.de/campus/2022-08/corona-nothilfen-bildungsministerium-studierende-dokumente}}.
%{https://www.zeit.de/campus/2022-08/corona-nothilfen-bildungsministerium-studierende-dokumente}}.

Durch die, auch aufgrund des russischen Angriffskrieges, nach wie vor steigenden Lebenshaltungskosten kämpfen Studierende im Alltag noch immer mit finanziellen Problemen. Auch eine Erhöhung des Mindestlohns konnte diese Entwicklung nicht aufhalten. Entsprechend ist für alle Studierenden ein finanzieller Ausgleich erforderlich!

Die bisherigen Anstrengungen der Bundesregierung, die Folgen dieser Preis\-stei\-ge\-run\-gen auszugleichen, sehen wir als nicht ausreichend an. Insbesondere die Einmalzahlung von \EUR{200} kam zu spät, war absolut unzureichend und unnötig kompliziert zu beantragen. Bei aktuell teilweise zweistelligen Inflationsraten\footnote{\url{https://www.destatis.de/DE/Themen/Wirtschaft/Konjunkturindikatoren/Basisdaten/vpi041j.html}}
%{https://www.destatis.de/DE/Themen/Wirtschaft/Konjunkturindikatoren/Basisdaten/\\vpi041j.html}}
liegt eine Belastung vor, welche mit dieser einmaligen Zahlung nicht auszugleichen ist.


Aus diesen Gründen fordert die ZaPF die aktuelle Bundesregierung auf, sich mit den Nöten und Belangen der Studierenden zu befassen und eine nachhaltige finanzielle Entlastung zu schaffen, indem sie unter anderem die Versprechen aus dem Koalitionsvertrag der Ampelkoalition umsetzen.


% \section*{Begründung}
% erfolgt mündlich

\vspace{1cm} 

\vfill
\begin{flushright}
	Verabschiedet am 01. Mai 2023 \\
	auf der ZaPF in Berlin.
\end{flushright}

\end{document}
